\bookchapter{第十四章：密碼管理器專題}

\chapterintro{喺數碼時代，密碼就係你嘅門匙。但匙太多、太複雜，點算？密碼管理器就係你嘅「萬能匙管家」。}

\booksection{一、點解你需要密碼管理器？}

現代人平均有幾十甚至上百個網絡帳號。如果每個帳號都用相同嘅密碼，一旦一個網站被駭，所有帳號都危險。如果每個帳號都用唔同嘅複雜密碼，根本記唔住。

密碼管理器解決咗呢個困境：你淨係需要記住一個主密碼，佢幫你記住所有其他密碼。

\booksection{二、密碼管理器嘅工作原理}

\begin{enumerate}
  \item \textbf{加密存儲}：所有密碼都經過高強度加密存儲，淨係你可以解密。
  \item \textbf{主密碼}：你設定一個主密碼，用於解鎖密碼庫。主密碼本身唔存儲喺伺服器上。
  \item \textbf{自動填充}：喺網站或App登入介面，密碼管理器自動填充帳號密碼。
  \item \textbf{密碼生成}：可以生成隨機、複雜嘅密碼，提高安全性。
\end{enumerate}

\booksection{三、加密原理簡介}

瞭解密碼管理器嘅加密原理，可以幫你判斷一個服務商係咪值得信任。下面用通俗嘅語言解釋幾個核心概念。

\booksubsection{AES-256 加密}

AES（高級加密標準）係美國國家標準與技術研究院（NIST）認可嘅對稱加密算法，被銀行、軍方、政府廣泛使用。「256」指密鑰長度為 256 位——理論上需要超過宇宙年齡嘅時間先可以暴力破解。主流密碼管理器（1Password、Bitwarden、KeePass 等）都用 AES-256 加密你嘅密碼庫。

\booksubsection{零知識架構（Zero-Knowledge Architecture）}

「零知識」唔係話服務商乜都唔存，而係話\textbf{加密同解密都喺你嘅設備本地完成，伺服器淨係保存已加密嘅密文}。具體流程：

\begin{enumerate}
  \item 你喺用戶端輸入主密碼。
  \item 用戶端用主密碼派生出加密密鑰（呢一步喺你手機/電腦上完成）。
  \item 淨係加密後嘅密文先上傳到伺服器。
  \item 下次登入時，伺服器將密文下發，用戶端再次用主密碼解密。
\end{enumerate}

關鍵點：\textbf{主密碼本身從不離開你嘅設備}，伺服器上都冇。所以即使服務商被駭，黑客攞到嘅都淨係一堆無法解密嘅密文——除非佢哋可以猜中你嘅主密碼。

\begin{tcbinfo}{點解「零知識」意味住服務商都睇唔到你嘅密碼？}
因為解密密鑰由你嘅主密碼派生，而主密碼淨係存在你腦中、淨係喺你本地設備上使用。服務商嘅伺服器上既冇主密碼，都冇解密密鑰，淨係得密文。呢個就係點解 Bitwarden、1Password 等正規服務商都明確聲明「我哋都無法睇到你嘅密碼」——唔係佢哋唔想睇，係技術上做不到。

\end{tcbinfo}

\booksubsection{密鑰派生函數：PBKDF2 與 Argon2}

主密碼唔可以直接用作加密密鑰——佢需要經過「密鑰派生函數」處理。呢個函數嘅作用係\textbf{故意等計算變慢}，等黑客暴力破解嘅成本高到不可行。

\begin{itemize}
  \item \textbf{PBKDF2}：老牌標準，被 iCloud 鑰匙串、Bitwarden、LastPass 等廣泛使用。通過反覆迭代（通常 10 萬次以上）雜湊嚟拖慢破解。
  \item \textbf{Argon2}：2015 年密碼雜湊競賽冠軍，比 PBKDF2 更現代。除咗「算力消耗」仲引入咗「記憶體消耗」——等 GPU/ASIC 等專用硬件都難以加速破解。Bitwarden、KeePassXC 等已支援。
\end{itemize}

迭代次數越高越安全，但解鎖速度都會變慢。正規服務商都會公開佢使用嘅派生算法同迭代次數，方便安全社區審計。

\booksubsection{案例警示：LastPass 2022 年數據洩露}

2022 年 8 月，老牌密碼管理器 LastPass 公告佢嘅開發環境被入侵；同年 11 月進一步披露，黑客攞走咗用戶嘅密碼庫備份（加密後嘅密文）。事件嘅關鍵教訓：

\begin{itemize}
  \item \textbf{零知識架構守住咗底線}：由於密碼庫本身係加密嘅，黑客攞到密文但解唔開——前提係用戶嘅主密碼足夠強。
  \item \textbf{弱主密碼用戶最危險}：黑客可以離線暴力破解弱主密碼，因此使用弱主密碼或重複主密碼嘅用戶面臨實質風險。
  \item \textbf{元數據都洩露}：雖然密碼內容加密，但站點 URL、備註等元數據可能未加密，會暴露用戶使用咗邊啲服務。
\end{itemize}

\begin{tcbwarn}{由 LastPass 事件學到嘅}
第一，主密碼必須長而且獨特（建議 16 位以上，包含大小寫、數字、符號）；第二，優先選擇支援 Argon2id、迭代次數高嘅服務商；第三，對高價值帳號（銀行、郵箱）務必開啟獨立嘅二次驗證，唔好淨係依賴密碼管理器。

\end{tcbwarn}

\booksection{四、主流密碼管理器推薦}

\booksubsection{系統自帶：}

\begin{itemize}
  \item \textbf{iCloud鑰匙串}：蘋果生態用戶首選。同Safari、App無縫整合，免費，安全。缺點係跨平臺使用不便。
  \item \textbf{Google密碼管理器}：安卓用戶方便。同Chrome、安卓系統整合，免費。
  \item \textbf{華為密碼保險箱}：鴻蒙系統內置，適合華為用戶。
\end{itemize}

\booksubsection{第三方專業工具：}

\begin{itemize}
  \item \textbf{Bitwarden}：開源免費，功能強大，支援所有平臺。強烈推薦畀注重私隱嘅用戶。
  \item \textbf{1Password}：介面美觀，功能豐富，適閤家庭同團隊。付費。
  \item \textbf{LastPass}：老牌密碼管理器，免費版功能受限。
  \item \textbf{KeePass}：完全本地存儲，開源免費。適合極客用戶。
\end{itemize}

\booksection{五、密碼強度真實評估}

密碼管理器普遍會顯示「密碼強度」指示條（紅/黃/綠），但呢個評估往往淨係睇長度同字符種類，並不全面。真正判斷密碼係咪安全，需要考慮「熵」同「係咪已洩露」。

\booksubsection{熵（Entropy）：密碼強度嘅科學度量}

「熵」衡量密碼嘅不確定性，單位係比特。熵越高，暴力破解所需嘗試次數越多。

\begin{itemize}
  \item \textbf{"123456"}：熵極低，瞬間破解。
  \item \textbf{"P@ssw0rd"}：睇落複雜，但出現喺所有字典攻擊列表入面，熵依然低。
  \item \textbf{"correct horse battery staple"}：四個隨機單詞，熵約 44 比特，破解難度遠超「P@ssw0rd」。
  \item \textbf{16 位隨機字符（含大小寫+數字+符號）}：熵約 100 比特，目前算力無法暴力破解。
\end{itemize}

結論：\textbf{「長度 + 隨機性」比「複雜符號」更重要}。呢個都係密碼管理器生成嘅隨機長密碼真正安全嘅原因。

\booksubsection{HaveIBeenPwned：查你嘅密碼係咪已洩露}

安全研究員 Troy Hunt 維護嘅\href{https://haveibeenpwned.com}{haveibeenpwned.com}收錄咗歷次大型數據洩露中流出嘅帳號同密碼。你可以：

\begin{enumerate}
  \item 輸入郵箱，查你嘅帳號係咪喺某次洩露中被曝光。
  \item 用 Pwned Passwords 功能查詢某個密碼係咪喺已知洩露庫中（密碼管理器同瀏覽器都已內置此查詢）。
\end{enumerate}

如果一個密碼出現喺洩露庫中，\textbf{無論佢睇落幾複雜都必須立即更換}——因為字典攻擊會優先嘗試呢啲已知密碼。

\begin{tcbinfo}{k-Anonymity 查詢原理}
擔心喺 HaveIBeenPwned 輸入密碼會再次洩露？唔使擔心。現代查詢採用「k-匿名」協議：瀏覽器淨係將密碼雜湊嘅前 5 位發畀伺服器，伺服器返所有以呢 5 位開頭嘅雜湊列表，本地再比對。即係話，伺服器永遠睇唔到你嘅完整密碼雜湊。1Password、Bitwarden 嘅「已洩露密碼檢測」都係咁樣工作嘅。

\end{tcbinfo}

\booksubsection{密碼管理器內置嘅監控功能}

\begin{itemize}
  \item \textbf{1Password Watchtower}：實時監控弱密碼、重複密碼、已洩露密碼、過期密碼、唔使用 2FA 嘅網站，並以紅/橙/黃分級展示。
  \item \textbf{Bitwarden Vault Health}：免費版即可查看弱密碼、重複密碼；付費版可檢測已洩露密碼（基於 HaveIBeenPwned）。
  \item \textbf{iCloud 鑰匙串「安全建議」}：iOS 設定 密碼 安全建議，會提示你邊啲密碼出現喺已知洩露中、邊啲重複使用。
  \item \textbf{Google 密碼管理器「密碼檢查工具」}：Chrome 內置，可以一鍵檢查已洩露、弱、重複密碼。
\end{itemize}

\begin{tcbpractice}{而家就查：你嘅密碼係咪安全}
\begin{enumerate}
  \item 打開 haveibeenpwned.com，輸入你嘅常用郵箱，睇下出現喺邊啲洩露事件中。
  \item iPhone 用戶：設定 密碼 安全建議，逐條處理紅色警告。
  \item Chrome 用戶：設定 密碼 檢查密碼，按提示修改。
  \item 1Password/Bitwarden 用戶：打開 Watchtower/Health Reports，優先處理「已洩露」同「重複」項。
\end{enumerate}

\end{tcbpractice}

\booksection{六、Passkey：無密碼登入嘅未來}

2024-2026 年，密碼技術正在經歷一場革命——\textbf{Passkey（通行密鑰）}有望徹底淘汰傳統密碼。Apple、Google、Microsoft 三大平臺聯合推動，FIDO 聯盟同萬維網聯盟（W3C）背書，呢個係數十年嚟登入方式最大嘅一次變革。

\booksubsection{乜嘢係 Passkey？}

Passkey 基於 FIDO2 / WebAuthn 標準，用\textbf{設備上嘅密鑰對}替代密碼。原理可以簡化為：

\begin{itemize}
  \item 你嘅手機生成一對密鑰：\textbf{公鑰}（註冊畀網站）+\textbf{私鑰}（永遠存在你設備嘅安全晶片中）。
  \item 登入時，網站發一個「挑戰」，你嘅手機用私鑰簽名，再由網站用公鑰驗證。
  \item 簽名前需要你通過生物識別（Face ID / 指紋）或設備 PIN 確認身份。
  \item 整個過程中，\textbf{冇任何密碼喺網絡中傳輸}。
\end{itemize}

\booksubsection{Passkey vs 傳統密碼}

\begin{center}
\begin{tabularx}{\linewidth}{|X|X|X|}
\hline
\textbf{維度} & \textbf{傳統密碼} & \textbf{Passkey} \tabularnewline
\hline
抗釣魚 & 差，釣魚網站可竊取密碼 & 天然抗釣魚，私鑰淨係對真實域名簽名 \tabularnewline
伺服器洩露影響 & 密碼雜湊可能被破解 & 伺服器淨係存公鑰，洩露無影響 \tabularnewline
便捷性 & 需記憶/輸入/管理 & Face ID/指紋一秒登入 \tabularnewline
重放攻擊 & 可能被中間人截獲重放 & 每次簽名一次性，無法重放 \tabularnewline
\hline
\end{tabularx}
\end{center}

\booksubsection{各平臺支援情況}

\begin{itemize}
  \item \textbf{Apple iCloud 鑰匙串 Passkey}：iOS 16+ / macOS Ventura+ 原生支援，Passkey 可以喺所有登入同一 Apple ID 嘅設備間同步，新設備登入網站時唔使重新創建。
  \item \textbf{Google 密碼管理器 Passkey}：安卓 + Chrome 原生支援，Passkey 存喺 Google 賬戶中，跨安卓設備同步；2023 年起 Google 個人帳號已支援用 Passkey 登入。
  \item \textbf{Microsoft}：Windows 11 / Edge 已支援 Passkey；微軟帳號、Outlook、Office 等服務陸續開放 Passkey 登入，並可以喺 Microsoft Authenticator 中管理。
  \item \textbf{1Password / Bitwarden}：已支援 Passkey 嘅創建、存儲、跨平臺同步，等你唔受限於單一生態。
\end{itemize}

\begin{tcbinfo}{跨設備同步係關鍵突破}
早期 Passkey 嘅最大痛點係「綁定單台設備」——換手機就要重新創建。2023 年起 Apple、Google、Microsoft 實現咗 Passkey 嘅雲端同步：iCloud 鑰匙串同步跨 Apple 設備，Google 密碼管理器同步跨安卓設備。1Password 等第三方就打通咗跨生態同步。呢個意味住 Passkey 已具備替代密碼嘅實用性。

\end{tcbinfo}

\booksubsection{點樣喺手機上創建同使用 Passkey}

\textbf{iOS（iPhone）路徑：}

\begin{enumerate}
  \item 喺支援 Passkey 嘅網站（如 Google、微軟、亞馬遜、eBay、PayPal）帳號安全頁選擇「添加通行密鑰 / Add passkey」。
  \item 系統彈出 Face ID / Touch ID 驗證視窗，通過後 Passkey 自動生成並儲存到 iCloud 鑰匙串。
  \item 下次登入該網站時，輸入用戶名後選擇「使用 iPhone 登入」，Face ID 一嘟即可。
  \item 管理入口：設定 密碼 揾到對應網站，可見 Passkey 條目。
\end{enumerate}

\textbf{安卓路徑：}

\begin{enumerate}
  \item 喺支援 Passkey 嘅網站帳號安全頁選擇「添加通行密鑰」。
  \item Google 密碼管理器彈窗 選擇你嘅 Google 帳號 螢幕鎖驗證（指紋/圖案/PIN）。
  \item Passkey 自動儲存到 Google 密碼管理器，跨同帳號安卓設備同步。
  \item 下次登入時，選擇「使用 Passkey」，螢幕鎖一解即可。
  \item 管理入口：安卓設定 密碼、密鑰同自動填充 密碼管理器 google.com。
\end{enumerate}

\booksubsection{Passkey 跨設備登入}

如果你喺 Windows 電腦上登入一個淨係得 iPhone 上有 Passkey 嘅網站，可以掃碼授權：

\begin{enumerate}
  \item 網站登入頁選擇「使用 Passkey」。
  \item 選擇「使用其他設備」。
  \item 電腦顯示二維碼，用 iPhone 掃碼。
  \item iPhone 彈出 Face ID 驗證，確認後網站喺電腦上完成登入。
\end{enumerate}

呢個過程通過藍牙+互聯網雙重加密通信，Passkey 私鑰唔會離開你嘅 iPhone。

\begin{tcbpractice}{而家就試：體驗無密碼登入}
\begin{enumerate}
  \item 登入 Google 帳號 安全 通行密鑰 添加通行密鑰，體驗 Face ID/指紋登入。
  \item 登入微軟帳號 安全 進階安全選項 添加通行密鑰。
  \item 喺淘寶、京東、eBay、PayPal 等已支援 Passkey 嘅服務上嘗鮮。
  \item 保留傳統密碼作為後備，逐步遷移高價值帳號到 Passkey。
\end{enumerate}

\end{tcbpractice}

\begin{tcbwarn}{過渡期建議}
Passkey 並未完全取代密碼——好多網站仍處於「Passkey + 密碼」並存階段。建議：高價值帳號優先開啟 Passkey，但\textbf{唔好立即刪除傳統密碼}，留作備用；同時繼續使用密碼管理器，因為唔係所有網站都支援 Passkey。將 Passkey 睇作「密碼管理器嘅強力補充」而唔係「立刻替換」。

\end{tcbwarn}

\booksection{七、硬件安全密鑰}

如果你對安全嘅要求高過普通用戶——比如持有大量加密貨幣、係高管/記者/活動人士、或保護企業核心帳號——硬件安全密鑰係比短訊驗證碼更可靠嘅選擇。

\booksubsection{乜嘢係硬件安全密鑰？}

硬件密鑰係一個物理設備（形似 U 盤），內部安全晶片存儲你嘅私鑰，需要物理接觸或 NFC 先可以完成驗證。代表產品：

\begin{itemize}
  \item \textbf{YubiKey}（Yubico 公司）：最主流，支援 FIDO2/WebAuthn、U2F、OTP、PIV、OpenPGP，USB-C / USB-A / NFC 等多種介面。5 系列 50-70 美元。
  \item \textbf{Google Titan}：Google 自家出品，USB-C + NFC / 藍牙，約 30-40 美元。
  \item \textbf{OnlyKey}：開源替代品，支援多種協議。
  \item \textbf{國產替代}：如飛天誠信 ePasskey，價格更親民。
\end{itemize}

\booksubsection{點解硬件密鑰比短訊驗證碼安全？}

\begin{itemize}
  \item \textbf{抗釣魚}：同 Passkey 一樣，硬件密鑰淨係對真實域名響應，釣魚網站無法騙取驗證碼。
  \item \textbf{抗 SIM 劫持}：短訊驗證碼可被 SIM-swap 攻擊截獲，硬件密鑰唔使手機號。
  \item \textbf{抗中間人}：FIDO2 協議每次簽名都包含挑戰值，無法重放。
  \item \textbf{物理隔離}：私鑰永遠喺硬件晶片內，惡意軟件即使攻陷你嘅電腦都讀唔到。
\end{itemize}

\booksubsection{適用場景}

\begin{itemize}
  \item 加密貨幣交易所帳號、錢包恢復。
  \item 企業 SSO、遠端辦公 VPN、SSH 登入伺服器。
  \item 郵箱、GitHub、社交媒體等「主帳號」——一旦被攻破可重置所有其他帳號。
  \item 記者、律師、高管等高敏感人羣。
\end{itemize}

\booksubsection{點樣配置硬件密鑰作為二次驗證}

以 Google 帳號 + YubiKey 為例：

\begin{enumerate}
  \item 購買 YubiKey（建議一次買兩把，一把日常用、一把備用，存放喺唔同地點）。
  \item 登入 myaccount.google.com 安全 兩步驗證。
  \item 選擇「添加安全密鑰」 將 YubiKey 插入 USB-C 口或貼近手機 NFC 感應區。
  \item 按提示觸摸 YubiKey 金屬觸點完成註冊。
  \item 建議註冊第二把備用密鑰。
  \item 喺「兩步驗證」設定入面，將硬件密鑰設為首選驗證方式，短訊驗證碼作為備用。
\end{enumerate}

主流支援嘅網站：Google、GitHub、Microsoft、Apple ID、Facebook、Twitter、Dropbox、GitLab、AWS、Coinbase、Binance 等。

\begin{tcbinfo}{雙密鑰策略}
硬件密鑰最大嘅風險係「遺失」。建議：(1) 同時註冊兩把唔同嘅密鑰；(2) 一把隨身攜帶（鑰匙扣），一把存放喺屋企保險櫃或銀行保管箱；(3) 部分網站允許將「備用恢復碼」列印出嚟離線保存。咁樣即使主密鑰遺失，都唔會被永久鎖喺帳號外。

\end{tcbinfo}

\begin{tcbwarn}{購買渠道要謹慎}
務必由官方店鋪或授權代理購買，唔好買二手或來源不明嘅密鑰——理論上預置過嘅硬件可能已被植入後門。YubiKey 嘅韌體喺出廠後無法更新，但 Yubico 會持續推送安全公告，關注佢官網即可。

\end{tcbwarn}

\booksection{八、家庭/團隊密碼共享}

家庭場景下，Netflix、家庭 Wi-Fi、路由器後台、智能門鎖密碼需要全家共用；團隊場景下，伺服器、數據庫、社交媒體帳號需要多人協作。\textbf{用微信發密碼、寫喺共享文檔入面都係高危做法}——一旦聊天記錄或文檔權限洩露，密碼就洩露咗。

\booksubsection{1Password 嘅「共享保險庫」（Vault）}

1Password 係家庭同團隊共享嘅標杆產品。佢嘅核心機制係\textbf{多個保險庫}：

\begin{itemize}
  \item \textbf{個人保險庫}：淨係你可以睇，如銀行、郵箱等私人帳號。
  \item \textbf{共享保險庫}：家人或團隊成員都可以睇，如 Netflix、家庭 Wi-Fi、共享信用卡。
  \item \textbf{權限分級}：團隊版可設定某些成員「淨係查看」或「淨係查看但不編輯」。
\end{itemize}

家人或團隊成員唔使共享主密碼，各自用自己嘅帳號登入，但可以睇到被授權嘅共享保險庫內容。密碼修改後所有成員實時同步。

\booksubsection{Bitwarden 嘅「組織」功能}

Bitwarden 嘅免費版支援個人密碼共享（通過「發送 Send」功能限時分享單條密碼）；付費版可創建「組織」，將若干條目共享畀家庭成員或團隊同事，權限可細分。Bitwarden 仲提供自託管選項，適合對數據主權要求高嘅企業。

\booksubsection{家庭場景：共享 Netflix、Wi-Fi 等密碼嘅安全方式}

\begin{itemize}
  \item \textbf{Netflix / 影片網站}：用 1Password 家庭版嘅「共享保險庫」存登入密碼，家庭成員各自登入即可，唔使喺羣入面發。
  \item \textbf{家庭 Wi-Fi 密碼}：iOS 設備間可通過 AirDrop 直接共享 Wi-Fi 密碼畀來訪嘅客人，唔使讀出嚟；安卓 12+ 都有「附近分享 Wi-Fi」功能。
  \item \textbf{智能門鎖、智能家電}：盡量用廠商 App 嘅「家庭成員」功能為每人單獨創建帳號，而唔係共享一個管理員帳號。
  \item \textbf{家用路由器後台}：管理員密碼存到 1Password 共享保險庫，訪客 Wi-Fi 單獨開一個 SSID。
\end{itemize}

\begin{tcbinfo}{1Password 家庭版嘅性價比}
1Password Families 約 4.99 美元/月，支援最多 5 位家庭成員，每人都有獨立嘅私人保險庫 + 4 個共享保險庫。計落人均唔到 1 美元/月，係目前家庭共享體驗最好嘅方案。Bitwarden 嘅家庭版 10 美元/年支援 6 人，性價比更高，但介面體驗略遜。

\end{tcbinfo}

\booksubsection{團隊場景：伺服器/數據庫密碼共享}

\begin{itemize}
  \item \textbf{1Password Teams / Business}：可按部門、項目組創建多個保險庫，細粒度權限控制，所有訪問行為可審計。
  \item \textbf{Bitwarden Teams / Enterprise}：開源方案，可自託管，企業版支援 SSO 整合、目錄同步。
  \item \textbf{職責分離原則}：管理員密碼唔應該淨係得一人掌握——設定 2-3 名「緊急聯絡人」，確保關鍵人員離職/休假時業務不中斷。
  \item \textbf{離職流程}：員工離職時立即撤銷其密碼管理器訪問權限，必要時旋轉所有佢接觸過嘅高敏感密碼。
\end{itemize}

\booksubsection{緊急訪問（Emergency Access）：萬一你出事}

想像一個場景：你獨自管理家庭所有重要帳號，但意外住院或不幸離世，家人無法登入銀行、郵箱、保險帳號。\textbf{「緊急訪問」功能}正係為此設計——授權一位信任嘅家人/朋友，喺你無法響應時獲取你嘅密碼庫。

\textbf{Bitwarden 緊急訪問設定：}

\begin{enumerate}
  \item 登入 Bitwarden Web Vault 設定 緊急訪問。
  \item 添加信任聯絡人嘅郵箱，邀請對方接受。
  \item 設定「等待期」（建議 7-14 天）：由對方發起請求到你確認嘅時間窗口。
  \item 如果你喺等待期內唔響應（拒絕或唔操作），到期後對方自動獲得訪問權。
  \item 你隨時可以手動拒絕請求或撤銷緊急訪問權限。
\end{enumerate}

\textbf{1Password 緊急訪問：}

\begin{enumerate}
  \item 1Password 唔直接叫「緊急訪問」，但通過「家庭/團隊組織」功能可實現——將家人加入家庭版後，將部分條目放入「共享保險庫」。
  \item 對個人版用戶，建議將主密碼告知一位高度信任嘅家人，或寫入密封信封存放於保險櫃。
  \item 1Password 仲提供「恢復碼」，開啟雙重認證時務必列印保存。
\end{enumerate}

\begin{tcbpractice}{而家就設：為家人留一把「備用匙」}
\begin{enumerate}
  \item 選定一位信任嘅家人作為「緊急聯絡人」。
  \item 喺密碼管理器中設定緊急訪問（Bitwarden）或加入家庭共享保險庫（1Password）。
  \item 將主密碼 + 恢復碼寫喺紙上，封入信封，存放於屋企保險櫃或告知可信家人。
  \item 每半年回顧一次，確認緊急聯絡人係咪仍然合適。
\end{enumerate}

\end{tcbpractice}

\booksection{九、密碼管理器使用最佳實踐}

\begin{itemize}
  \item \textbf{設定強主密碼}：主密碼係最後一道防線，務必設定一個長而且複雜嘅密碼，並牢記。
  \item \textbf{啟用雙重認證}：為密碼管理器帳號開啟雙重認證，增加一層保護。
  \item \textbf{定期檢查弱密碼}：大多數密碼管理器有「安全審計」功能，幫你揾出弱密碼同重複密碼。
  \item \textbf{唔好忘記主密碼}：主密碼無法揾返，忘記就意味住遺失所有密碼。建議寫喺紙上，放喺安全嘅地方。
  \item \textbf{緊急訪問}：設定緊急聯絡人，萬一你無法訪問，信任嘅人可以獲取你嘅密碼。
\end{itemize}

\begin{tcbwarn}{重要提醒}
密碼管理器雖然方便，但都有風險：如果主密碼洩露，所有密碼都暴露咗。所以主密碼絕對唔可以同其他密碼相同，都唔可以寫喺容易被揾到嘅地方。另外，建議選擇有良好安全記錄嘅服務商。

\end{tcbwarn}

\begin{tcbpractice}{而家就試：開始使用密碼管理器}
\begin{enumerate}
  \item 選擇一個密碼管理器（推薦Bitwarden或系統自帶）。
  \item 註冊帳號，設定一個強主密碼。
  \item 匯入現有密碼（如果有嘅話）或開始添加新密碼。
  \item 安裝瀏覽器插件或手機App，體驗自動填充功能。
  \item 檢查安全審計報告，揾出需要改進嘅密碼。
\end{enumerate}

\end{tcbpractice}
